\SourcePage{121}
\section{Fukaya 范畴：初步版本}\label{sec:9}
\subsection{族}\label{sec:9a}
\begin{definition*}
设 $\widehat S$ 是某个光滑定向带边曲面，$\Sigma=\Sigma^+\sqcup\Sigma^-$ 是有限个边界点所成的集合，$S=\widehat S\setminus\Sigma$。取纤维丛
\[
\pi:\mathcal S\longrightarrow\mathcal R,
\]
其纤维为 $S$，结构群是 $\widehat S$ 上在 $\Sigma$ 处等于恒等映射的保定向微分同胚群。每个纤维 $\mathcal S_r$ 都有自然紧化 $\widehat{\mathcal S}_r$，它们典范地组成光滑固有纤维丛 $\widehat\pi:\widehat{\mathcal S}\to\mathcal R$。此外，$\Sigma^\pm$ 中的每个点都给出截面 $\zeta:\mathcal R\to\partial\widehat{\mathcal S}$，满足 $\mathcal S=\widehat{\mathcal S}\setminus\bigcup\zeta(\mathcal R)$。常将 $\Sigma^\pm$ 中的点与对应的截面等同。现在假设，$\mathcal S$ 的纤维方向切丛上配备了殆复结构 $I_{\mathcal S}$，并能延伸到 $\widehat{\mathcal S}$，从而使每个纤维 $\mathcal S_r$ 成为边界带标点的 Riemann 曲面。在这种情形，称 $\mathcal S$ 为\emph{边界带标点的曲面族}。
\end{definition*}

$\mathcal S\to\mathcal R$ 的一组条带状端，由对每个 $\zeta\in\Sigma^\pm$ 给定的、覆盖 $\mathcal R$ 的固有嵌入 $\epsilon_\zeta:\mathcal R\times Z^\pm\to\mathcal S$ 组成，它们在各纤维上限制为条带状端。这些嵌入自动延伸为光滑、逐纤维全纯的嵌入 $\mathcal R\times\widehat Z^\pm\to\widehat{\mathcal S}$，其在 $\pm\infty\in\widehat Z^\pm$ 处的限制，就是对应截面 $\zeta$。任意族都容许条带状端，根本原因是：对单个边界带标点的 Riemann 曲面，端的选择空间可缩。这与在 Riemann 曲面内部一点附近选取局部坐标的情形不同，后者有本质性的 $S^1$ 自由度，并由此产生引力后裔。

换一个稍有不同的角度，假设已经固定一组端 $\{\epsilon_\zeta\}$。此时通常方便采用使端保持不变的 $\mathcal S\to\mathcal R$ 的局部平凡化，其含义如下。选择一点 $r_0\in\mathcal R$，设 $S=\mathcal S_{r_0}$ 为该点上的纤维，条带状端为 $\epsilon_{S,\zeta}=\epsilon_\zeta(r_0,\cdot):Z^\pm\to S$。对充分小的邻域 $r_0\in U\subset\mathcal R$，存在逐纤维微分同胚 $\Psi:U\times S\to\mathcal S|_U$，它在 $\{r_0\}\times S$ 上为恒等映射，并满足
\begin{equation}\label{eq:9.1}
\epsilon_\zeta(r,s,t)=\Psi(r,\epsilon_{S,\zeta}(s,t)),\qquad r\in U,\ (s,t)\in Z^\pm.
\end{equation}
于是，$\Psi^*I_{\mathcal S}$ 是由 $r\in U$ 参数化的 $S$ 上复结构族，只在紧致子集 $S\setminus\bigcup\epsilon_\zeta(Z^\pm)$ 上变化。

\subsection{Kodaira–Spencer 映射}\label{sec:9b}
设 $S$ 是带有条带状端的 Riemann 曲面。考虑切丛上的标准 $\overline\partial$ 算子，边界条件取在全实子丛 $T(\partial S)$ 中。按\secref{8h}的术语，它是容许的，但不是非退化的，因此在带有充分小权重 $\epsilon>0$ 的加权 Sobolev 空间上成为 Fredholm 算子，见\eref{8.16}。实际上，希望容许稍大一些的定义域
\begin{equation}\label{eq:9.2}
W^{k,p;\epsilon}(S,TS,T(\partial S))^+\supset W^{k,p;\epsilon}(S,TS,T(\partial S)),
\end{equation}
\SourcePage{122}
由在每个条带状端上具有形式 $\epsilon_\zeta^*Z=c_\zeta\partial_s+\widetilde Z_\zeta$ 的向量场 $Z$ 组成，其中 $c_\zeta\in\R$、$\widetilde Z_\zeta\in W^{k,p;\epsilon}$。这只增加有限维，每个端增加一维，对应于常数 $c_\zeta$；因此 $\overline\partial$ 仍可延伸为 Fredholm 算子
\begin{equation}\label{eq:9.3}
\overline\partial^+:W^{k,p;\epsilon}(S,TS,T(\partial S))^+
\longrightarrow W^{k-1,p;\epsilon}(S,\Omega_S^{0,1}\otimes TS).
\end{equation}
将其余核记为 $\mathcal T_S$；它既不依赖于 $(k,p)$，也不依赖于充分小的 $\epsilon>0$。若 $S$ 稳定，即不存在连续的自同构族，则\eref{9.3}是单射，且 $\dim\mathcal T_S=|\Sigma|-3\chi(\widehat S)$。\znote{114}{按\eref{9.3}，$\mathcal T_S$ 是 $\overline\partial^+$ 的余核，可看作 $S$ 的模空间的切空间；当 $S$ 是带 $d+1$ 个标点的圆盘时 $|\Sigma|=d+1$、$\chi(\widehat S)=1$，故 $\dim\mathcal T_S=d-2$，与 $\mathcal R^{d+1}$ 的维数相符，这正是 $d\geq2$ 时模空间非空且无多余参数的维数解释。单射性用到 $S$ 稳定：若存在连续自同构族，$\overline\partial^+$ 会多出核，$\mathcal T_S$ 就不再是切空间。[来源：原书 §(9b) 末尾所引 [72]（带标点曲线的代数变形理论），第 94 页；维数计算并见 §(9c) 对 $\mathcal R^{d+1}$ 的定义。]}

设 $\mathcal S\to\mathcal R$ 是带有条带状端的 Riemann 曲面族。选择一点 $r_0\in\mathcal R$，将其纤维记为 $S=\mathcal S_{r_0}$，并取一个像\eref{9.1}那样使条带状端保持不变的局部平凡化 $\Psi$。对 $S$ 上拉回复结构族 $(\Psi^*I_{\mathcal S})$ 求导，便为 $(T\mathcal R)_{r_0}$ 中的每个切向量给出
\[
C_c^\infty(S,\End^{0,1}(TS))=C_c^\infty(S,\Omega_S^{0,1}\otimes TS)
\]
中的一个元素。向 $\operatorname{coker}(\overline\partial^+)$ 投影，定义 \emph{Kodaira–Spencer 映射}
\begin{equation}\label{eq:9.4}
\kappa_{r_0}:(T\mathcal R)_{r_0}\longrightarrow\mathcal T_S.
\end{equation}
它不依赖于平凡化，也不依赖于条带状端的选择。后一事实从定义中并不立即可见，因为我们偏好采用与条带状端结构紧密相关的函数空间。
\begin{definition*}
若对每个 $r_0\in\mathcal R$，\eref{9.4}都是双射，就称 $\mathcal S\to\mathcal R$ 为\emph{无穷小极小通用族}。
\end{definition*}
为了说明这一名称确有所据，应该证明这样的族具有局部通用性质。这并不困难，但我们实际上用不到它。有兴趣的读者可参看带标点 Riemann 曲面的代数几何变形理论\bcite{72}，第 94 页；两者完全类似。

\subsection{带标点圆盘}\label{sec:9c}
现在将讨论限制于与 Fukaya 范畴相关的一类曲面。
\begin{definition*}
对于 $d\geq0$，一个\emph{带 $(d+1)$ 个标点的圆盘}，是紧化 $\widehat S\cong D$ 为圆盘、并有一个入向和 $d$ 个出向无穷远点的边界带标点 Riemann 曲面 $S$。约定按沿圆盘边界的循环次序给无穷远点编号，使入向点为 $\Sigma^-=\{\zeta_0\}$，出向点为 $\Sigma^+=\{\zeta_1,\ldots,\zeta_d\}$；若 $S$ 有条带状端，则相应编号为 $\epsilon_0,\ldots,\epsilon_d$。$\partial S$ 的连通分支按循环次序记为 $C_0,\ldots,C_d$，其中 $\zeta_0$ 位于 $C_d$ 与 $C_0$ 之间；$d=2$ 的情形见\sr{fig:9.1}{图~9.1}。若 $S$ 带有 Lagrangian 标记，则相应记为 $L_0,\ldots,L_d$。与无穷远点对应的子流形对于是为
\begin{equation}\label{eq:9.5}
(L_{\zeta_k,0},L_{\zeta_k,1})=\begin{cases}(L_0,L_d),&k=0,\\(L_{k-1},L_k),&k>0.\end{cases}
\end{equation}

\end{definition*}
\SourcePage{123}
\SourceFigure{figure-9-1}{0.30946121}{9.1}
在稳定范围 $d\geq2$ 内，带 $(d+1)$ 个标点的圆盘没有非平凡自同构，并存在它们的万有族，记为
\begin{equation}\label{eq:9.6}
\mathcal S^{d+1}\longrightarrow\mathcal R^{d+1}.
\end{equation}
构造很直接。设 $\operatorname{Conf}_{d+1}(\partial D)\subset(\partial D)^{d+1}$ 是圆周上编号与循环次序相容的 $(d+1)$ 元点组所成的构形空间。$\operatorname{Aut}(D)\cong\operatorname{PSL}_2(\R)$ 自由且固有地作用在此空间上，令
\begin{equation}\label{eq:9.7}
\begin{aligned}
\mathcal R^{d+1}&=\operatorname{Conf}_{d+1}(\partial D)/\operatorname{Aut}(D),\\
\widehat{\mathcal S}^{d+1}&=\operatorname{Conf}_{d+1}(\partial D)\times_{\operatorname{Aut}(D)}D.
\end{aligned}
\end{equation}
投影 $\widehat{\mathcal S}^{d+1}\to\mathcal R^{d+1}$ 有截面 $\zeta_k([z_0,\ldots,z_d])=[z_0,\ldots,z_d,z_k]$；按照一般步骤，令 $\mathcal S^{d+1}=\widehat{\mathcal S}^{d+1}\setminus\bigcup_k\zeta_k(\mathcal R^{d+1})$。参数化形式的 Riemann 单值化定理表明，\eref{9.6}确实是万有族：每个带 $(d+1)$ 个标点的圆盘族，都由从其底空间到 $\mathcal R^{d+1}$ 的唯一光滑映射拉回 $\mathcal S^{d+1}$ 而得。类似的无穷小结论也不难证明，它说明\eref{9.6}是无穷小极小通用的。
\begin{definition*}
\emph{条带状端的万有选择}，是对每个 $d\geq2$，在 $\mathcal S^{d+1}\to\mathcal R^{d+1}$ 上选取一组条带状端 $\{\epsilon_k^{d+1}\}_{0\leq k\leq d}$。当然，其万有性意指：任何带 $(d+1)$ 个标点的圆盘族，都通过其分类映射继承条带状端。
\end{definition*}

\clearpage
\subsection{树的术语}\label{sec:9d}
涉及多个带标点圆盘的粘合运算，可以自然地用树来描述。这里汇集有关术语，其中一些将在后面使用。
\begin{definition*}
对于 $d>0$，\emph{$d$ 叶树}是固有嵌入的平面树 $T\subset\R^2$，具有 $d+1$ 条半无限边：指定其中一条称为\emph{根}，其余称为\emph{叶}。半无限边也称为\emph{外边}，其他边称为\emph{内边}。以 $\operatorname{Ve}(T)$ 表示顶点集，以
\SourcePage{124}
$\operatorname{Ed}^{\mathrm{int}}(T)$ 表示内边集。回顾图的\emph{旗标}是由一个顶点与一条相邻边构成的二元组 $f=(v,e)$。用 $\operatorname{Fl}^{\mathrm{int}}(T)$ 表示\emph{内旗标}的集合，即其中的边为内边的旗标。

按从根向叶“向上”流动的方向给 $T$ 定向，见\sr{fig:9.2}{图~9.2}。若 $e$ 的定向背离 $v$，则称旗标 $f=(v,e)$ 为\emph{正}的，否则称为\emph{负}的。每条内边 $e$ 对应两个旗标 $f^\pm(e)$，一正一负。同样，每条外边有唯一的旗标；边为叶时它为正，边为根时它为负。有时方便将与给定顶点 $v$ 相邻的旗标编号为 $f_0(v),\ldots,f_{|v|-1}(v)$，其中 $|v|$ 是顶点的度数。编号总是从唯一的负旗标 $f_0(v)$ 开始，按给定平面嵌入的逆时针方向继续。
\end{definition*}
\SourceFigure{figure-9-2}{0.43615338}{9.2}
\begin{definition*}
若所有顶点 $v$ 都满足 $|v|\geq3$，则称 $d$ 叶树\emph{稳定}；若 $|v|\geq2$，则称其\emph{半稳定}。
\end{definition*}

\subsection{树与粘合}\label{sec:9e}
假设给定一棵 $d$ 叶树 $T$；对每个顶点 $v$，给定带有条带状端的 $|v|$ 标点圆盘 $S_v$；对每条内边 $e$，给定粘合长度 $l_e>0$。从这些数据出发，粘合构造产生一个新的带条带状端的 $(d+1)$ 标点圆盘 $S_l$。

先注意，旗标 $f=(v,e)$ 与 $S_v$ 中的无穷远点之间存在典范双射，使负、正旗标分别对应入向、出向点。具体说，将 $S_v$ 的第 $k$ 个无穷远点对应于旗标 $f_k(v)$。相应调整记号，以 $\epsilon_f$ 表示 $S_v$ 在对应于 $f$ 的无穷远点附近的条带状端。粘合过程的第一步，是从每个内旗标对应的端切掉一段：
\begin{equation}\label{eq:9.8}
S'_v=S_v\setminus\left(\coprod_{\substack{f=(v,e)\\f\in\operatorname{Fl}^{\mathrm{int}}(T)}}\epsilon_f(Z_{l_e}^\pm)\right),
\end{equation}
\SourcePage{125}
其中按旗标为负或为正，分别采用 $Z_{l_e}^-=(-\infty,-l_e)\times[0,1]$ 或 $Z_{l_e}^+=(l_e,\infty)\times[0,1]$。然后，将这些截断端的剩余部分彼此识别：令 $S_l=\coprod_vS'_v/\sim$，其中对每条内边 $e$ 以及 $(s,t)\in[0,l_e]\times[0,1]$，规定
\begin{equation}\label{eq:9.9}
\epsilon_{f^-(e)}(s-l_e,t)\sim\epsilon_{f^+(e)}(s,t).
\end{equation}
$S_l$ 的条带状端继承自 $T$ 的外旗标所对应的端。更精确地，从根开始按循环次序将这些旗标编号为 $f_0,\ldots,f_d$，则 $S_l$ 的第 $k$ 个条带状端为
\begin{equation}\label{eq:9.10}
\epsilon_{l,k}=\epsilon_{f_k}:Z^\pm\longrightarrow\coprod_vS'_v\longrightarrow S_l.
\end{equation}
当 $T$ 只有一个顶点时，粘合过程没有作任何改变，只是返回原曲面；若有两个顶点，它就等于\eref{8.5}中的基本连通和运算；一般情况下，它可看作有限次这种运算的复合，执行次序任意。

\begin{note}\label{st:9.1}
曲面 $S_l$ 继承一个\emph{厚薄分解}。该分解来自粘合过程，并非曲面的内在结构，因此不同于双曲几何所给出的同名分解。薄部由条带状端及长度为 $l_e$ 的有限条带组成：
\[
\begin{aligned}
S_l^{\mathrm{thin}}&=\coprod_{k=0}^d\epsilon_{l,k}(Z^\pm)\ \sqcup\!
\coprod_{e\in\operatorname{Ed}^{\mathrm{int}}(T)}\epsilon_{f^+(e)}([0,l_e]\times[0,1]),\\
S_l^{\mathrm{thick}}&=\coprod_v\left(S_v\setminus\coprod_{k=0}^{|v|-1}\epsilon_{f_k(v)}(Z^\pm)\right).
\end{aligned}
\]

\end{note}
后面将看到，最好采用长度的负指数作为\emph{粘合参数}。更精确地，$\rho_e\in(-1,0)$ 与相应长度的关系为 $l_e=-\log(-\rho_e)/\pi$，因此 $\rho_e\to0$ 对应于 $l_e\to\infty$。粘合后的曲面 $S_\rho=S_l$ 组成一个 $(d+1)$ 标点圆盘族
\[
\mathcal S\longrightarrow\mathcal R=(-1,0)^{\operatorname{Ed}^{\mathrm{int}}(T)},
\]
其条带状端 $\mathcal R\times Z^\pm\to\mathcal S$ 在各纤维上由\eref{9.10}给出。从某种意义上说，也可容许某些粘合参数为零的极限情形。此时的规则是：在粘合过程的两个步骤\eref{9.8}、\eref{9.9}中，都略去退化边，即粘合参数为零的边。结果不再是单个带标点圆盘，而是若干这种圆盘的不交并；它由
\SourcePage{126}
将 $T$ 的所有非退化边压缩后得到的树 $\overline T$ 描述。特别地，$\overline T$ 的每个顶点 $\overline v$ 对应一个连通分支，它是带 $|\overline v|$ 个标点的圆盘。在所有粘合参数都为零的平凡情形，所得正是各 $S_v$ 的不交并。\znote{118}{此处“压缩”（收缩）内边的效果，是把该边两个端点所代表的标点合并成一个：粘合参数非零的边两侧恰被实际粘合，故在树上应把它收缩掉；参数为零的边不作粘合，曲面沿它分裂，但该边仍留在树上，用以记录各连通分支之间的关系。于是每个顶点恰好对应一个连通分支，顶点的度数就是该分支上的标点数。所有参数为零时树不变，曲面回到各 $S_v$ 的不交并，与正文最后一句一致。[来源：原书 §(9e) 的粘合构造，即\eref{9.8}、\eref{9.9}；本注只是该处逻辑的展开。]}将这些纤维加入 $\mathcal S$，得到部分紧化 $\overline{\mathcal S}$；它是光滑带角流形，并带有一个浸没映射 $\overline{\mathcal S}\to\overline{\mathcal R}=(-1,0]^{\operatorname{Ed}^{\mathrm{int}}(T)}$，显然这不是纤维化。

粘合构造的一种直接推广，是对每个顶点 $v$，从带有条带状端的 $|v|$ 标点圆盘族 $\mathcal S_v\to\mathcal R_v$ 出发。粘合得到新的族
\[
\mathcal S\longrightarrow\mathcal R=(-1,0)^{\operatorname{Ed}^{\mathrm{int}}(T)}\times\prod_v\mathcal R_v.
\]
像前面一样，容许粘合参数为零，便得到部分紧化 $\overline{\mathcal S}\to\overline{\mathcal R}$。

\subsection{Deligne–Mumford–Stasheff 紧化：结合多面体}\label{sec:9f}
作为集合，$\mathcal R^{d+1}$ 的 Deligne–Mumford–Stasheff 紧化为
\begin{equation}\label{eq:9.11}
\overline{\mathcal R}^{d+1}=\coprod_T\mathcal R^T.
\end{equation}
这里对所有稳定 $d$ 叶树 $T$ 取并，并令 $\mathcal R^T=\prod_{v\in\operatorname{Ve}(T)}\mathcal R^{|v|}$。为在\eref{9.11}上赋予拓扑与光滑结构，先作某个条带状端的万有选择，再采用粘合构造。具体地，取一棵稳定树 $T$，将万有族 $\mathcal S^{|v|}\to\mathcal R^{|v|}$ 配给每个顶点 $v$。将它们粘合，得到 $(d+1)$ 标点圆盘族
\begin{equation}\label{eq:9.12}
\mathcal S\longrightarrow\mathcal R=(-1,0)^{\operatorname{Ed}^{\mathrm{int}}(T)}\times\mathcal R^T,
\end{equation}
按定义，它由光滑的\emph{粘合映射}
\begin{equation}\label{eq:9.13}
\gamma^T:\mathcal R\longrightarrow\mathcal R^{d+1}
\end{equation}
分类。像前面所说明的，可以容许部分或全部粘合参数为零。结果是退化的粘合曲面，即由压缩后的树 $\overline T$ 控制的若干带标点圆盘的不交并。显然，这确定了\eref{9.11}中 $\overline T$-层的一点，从而定义粘合映射的典范延伸
\begin{equation}\label{eq:9.14}
\overline\gamma^T:\overline{\mathcal R}=(-1,0]^{\operatorname{Ed}^{\mathrm{int}}(T)}\times\mathcal R^T\longrightarrow\overline{\mathcal R}^{d+1}.
\end{equation}
在子集 $\{0\}^{\operatorname{Ed}^{\mathrm{int}}(T)}\times\mathcal R^T$ 上，它代表平凡的、退化程度最大的粘合运算，因此就是\eref{9.11}中 $T$-层上的恒等映射。
\SourcePage{127}
Deligne–Mumford–Stasheff 紧化的拓扑，由所有 $T$ 的映射\eref{9.14}诱导。也就是说，从 $\overline{\mathcal R}^{d+1}$ 到任意空间的映射连续，当且仅当它与所有 $\overline\gamma^T$ 的复合连续。可微结构由类似性质刻画，只需把“连续”替换为“光滑”。这样，$\overline{\mathcal R}^{d+1}$ 成为紧致带角流形，而且不依赖于条带状端的万有选择；此外，\eref{9.11}就是其内部及各连通边界层的分解。\znote{119}{当所有标点都在边界上时，$\overline{\mathcal R}^{d+1}$ 其实就是 Stasheff 结合多面体：它是 $d-2$ 维紧致带角流形，其面按标点的括号化方式编号，$K_3$ 是区间、$K_4$ 是五边形（两个二顶点树构成五边形的边，三个二顶点树构成其顶点，即\secref{9g}的五边形论证）。一族这样的多面体构成消解结合律的拓扑 operad，这正是\eref{9.11}的分层由树给出、且边界层按内边数降维的原因。[来源：nLab "associahedron"，https:/\allowbreak /\allowbreak ncatlab.org/\allowbreak nlab/\allowbreak show/\allowbreak associahedron（2026-11-15 访问）。]}

刚才陈述的性质并不能从构造中立即看出；事实上，甚至 $\overline{\mathcal R}^{d+1}$ 上的拓扑是否 Hausdorff，都可能还不清楚。验证这些性质的最佳方式，是采用 $\overline{\mathcal R}^{d+1}$ 的另一种定义，将它视为复 Deligne–Mumford 空间 $\overline{\mathcal M}_{0,d+1}$ 的实点集 $\overline{\mathcal M}_{0,d+1}(\R)$ 的一部分；这一解释见\bcite{58}、\bcite{60}。按这种方式引入时，$\overline{\mathcal R}^{d+1}$ 自带紧致光滑带角流形的典范结构。
\begin{lemma}\label{st:9.2}
在 $\overline{\mathcal R}^{d+1}$ 上赋予由其嵌入 $\overline{\mathcal M}_{0,d+1}(\R)$ 得到的光滑结构。则延伸后的粘合映射 $\overline\gamma^T$ 光滑，而且在所有粘合参数都为零的子集附近是局部微分同胚。
\end{lemma}
\begin{proof}[证明概要]
先回顾，嵌入 $\mathcal R^{d+1}\subset\overline{\mathcal M}_{0,d+1}(\R)\subset\overline{\mathcal M}_{0,d+1}$ 是通过将带标点圆盘“双倍化”为带有反全纯对合的穿孔 $\CP^1$ 实现的。圆盘上的条带状端也可双倍化，得到 $\CP^1$ 各穿孔附近的局部全纯坐标。将粘合参数复化，就能把\eref{9.14}延伸为映射
\begin{equation}\label{eq:9.15}
(\overline\gamma^T)^{\CC}:\operatorname{int}(D)^{\operatorname{Ed}^{\mathrm{int}}(T)}\times\mathcal R^T\longrightarrow\overline{\mathcal M}_{0,d+1},
\end{equation}
其中 $\operatorname{int}(D)=D\setminus\partial D$ 是开单位圆盘。更精确地，固定 $\mathcal R^T$ 中的一点，复化的粘合构造就在 $\operatorname{int}(D)^{\operatorname{Ed}^{\mathrm{int}}(T)}$ 上产生稳定带标点曲线的全纯族。\eref{9.15}是这个族的分类映射，因此它关于粘合参数全纯，并且在所有粘合参数都非零的开子集上光滑。

为推出它处处光滑，采用 Cauchy 公式的以下初等推论：设 $\overline U=\operatorname{int}(D)^q\times\R^p$、$U=(\operatorname{int}(D)\setminus\{0\})^q\times\R^p$，并对 $r\in\R^p$ 令 $\overline U_r=\operatorname{int}(D)^q\times\{r\}$。假设函数 $g:\overline U\to\CC$ 在 $U$ 上光滑，且在每个 $\overline U_r$ 上的限制都全纯，则 $g$ 光滑。最后，为证明 $\overline\gamma^T$ 是局部微分同胚，需要在所有粘合参数都为零的点处，计算其关于粘合参数的导数。复几何，更确切地说是带节点 Riemann 曲面的变形理论，再次使这一点成为简单的练习。
\end{proof}
有了\stref{引理}{9.2}，很容易看出，最初通过任意条带状端万有选择的粘合映射\eref{9.14}定义的 $\overline{\mathcal R}^{d+1}$ 的拓扑与光滑结构，
\SourcePage{128}
和嵌入 $\overline{\mathcal R}^{d+1}\subset\overline{\mathcal M}_{0,d+1}(\R)$ 给出的结构相同。这就推出了上面陈述的所有性质。

定义 $\mathcal S^T\to\mathcal R^T$ 为万有纤维化 $\mathcal S^{|v|}\to\mathcal R^{|v|}$（$v\in\operatorname{Ve}(T)$）的各拉回的不交并；$\mathcal S^T$ 的纤维是若干 $|v|$ 标点圆盘的不交并。万有族 $\mathcal S^{d+1}\to\mathcal R^{d+1}$ 有一个部分紧化，在集合意义下为
\begin{equation}\label{eq:9.16}
\overline{\mathcal S}^{d+1}=\coprod_T\mathcal S^T\longrightarrow\overline{\mathcal R}^{d+1}.
\end{equation}
通过族的总空间上的粘合映射，可在其中赋予拓扑及光滑结构。按定义，与\eref{9.12}对应的粘合映射 $\gamma^T$ 由逐纤维同构
\[
\Gamma^T:\mathcal S\longrightarrow\mathcal S^{d+1}
\]
覆盖。在解释一般粘合构造时已经指出，$\mathcal S$ 有自然的部分紧化 $\overline{\mathcal S}\to\overline{\mathcal R}$。将它与\eref{9.16}比较，可见 $\Gamma^T$ 延伸为 $\overline\Gamma^T:\overline{\mathcal S}\to\overline{\mathcal S}^{d+1}$。像前面一样，这些映射所诱导的拓扑与光滑结构，和把 $\overline{\mathcal S}^{d+1}$ 嵌入 $\overline{\mathcal M}_{0,d+1}$ 上的稳定曲线万有族得到的结构相同。
